\section{Model}

Physical modelling parameters and time-stepping configuration for the Lagrangian particle dispersion simulation.

\begin{lstlisting}[language=C++]
Model {
    <model parameters>
}
\end{lstlisting}

\subsection*{Parameters}

D
\vspace{\the\parameterListSpacing}
\begin{itemize}
    \item {\bf Description}: Molecular diffusion coefficient $D$ of the scalar quantity being dispersed.
    \item {\bf Required/Optional}: Required
    \item {\bf Type}: Real number
    \item {\bf Constraints}: Must be $\geq 0$
    \item {\bf Example}: \lstinline!D = 1.5e-5!
\end{itemize}

turbulentSchmidtNumber
\vspace{\the\parameterListSpacing}
\begin{itemize}
    \item {\bf Description}: Turbulent Schmidt number $Sc_t$, used to relate the turbulent eddy diffusivity to the eddy viscosity from the flow field. The turbulent eddy diffusivity is computed as $D_t = \nu_t / Sc_t$.
    \item {\bf Required/Optional}: Required
    \item {\bf Type}: Real number
    \item {\bf Constraints}: Must be $> 0$
    \item {\bf Example}: \lstinline!turbulentSchmidtNumber = 0.7!
\end{itemize}

numberOfTimeSteps
\vspace{\the\parameterListSpacing}
\begin{itemize}
    \item {\bf Description}: Total number of timesteps to simulate.
    \item {\bf Required/Optional}: Required
    \item {\bf Type}: Integer
    \item {\bf Constraints}: Must be $> 0$
    \item {\bf Example}: \lstinline!numberOfTimeSteps = 5000!
\end{itemize}

timeStepSize
\vspace{\the\parameterListSpacing}
\begin{itemize}
    \item {\bf Description}: Size of each timestep $\Delta t$. The total simulation time then becomes numberOfTimeSteps $\times \Delta t$ .
    \item {\bf Required/Optional}: Required
    \item {\bf Type}: Real number
    \item {\bf Constraints}: Must be $> 0$
    \item {\bf Example}: \lstinline!timeStepSize = 0.05!
\end{itemize}

particleSplitTimeStepInterval
\vspace{\the\parameterListSpacing}
\begin{itemize}
    \item {\bf Description}: Interval in timesteps at which particle splitting is performed. At each time interval, 
        particles are split into two, which each split particle getting half the mass of the parent. 
        Particle splitting subdivides particles to maintain numerical accuracy as the plume disperses.
    \item {\bf Required/Optional}: Required
    \item {\bf Type}: Integer
    \item {\bf Constraints}: Must be $> 0$
    \item {\bf Example}: \lstinline!particleSplitTimeStepInterval = 50!
\end{itemize}

maxNumberOfParticleSplits
\vspace{\the\parameterListSpacing}
\begin{itemize}
    \item {\bf Description}: Maximum number of times partile splitting is performed.
    \item {\bf Required/Optional}: Required
    \item {\bf Type}: Integer
    \item {\bf Constraints}: Must be $\geq 0$. A value of 0 disables splitting entirely.
    \item {\bf Example}: \lstinline!maxNumberOfParticleSplits = 4!
\end{itemize}

initialParticlesPerUnitMass
\vspace{\the\parameterListSpacing}
\begin{itemize}
    \item {\bf Description}: Number of computational particles used to represent one unit of released mass at the time of emission. Controls the initial resolution of the particle representation.
    \item {\bf Required/Optional}: Required
    \item {\bf Type}: Real number
    \item {\bf Constraints}: Must be $> 0$
    \item {\bf Example}: \lstinline!initialParticlesPerUnitMass = 100!
\end{itemize}
